\begin{abstract}
Bangladesh's electric easy-bikes run overwhelmingly on lead-acid batteries that last about a year, waste grid electricity and feed an informal recycling chain that is a major source of lead exposure. Lithium iron phosphate (LFP) and nickel-manganese-cobalt (NMC) packs cost more up front, and their advantage under Bangladesh's hot climate has not been quantified. We build a transparent, temperature-driven MATLAB/Simulink model that couples NASA POWER 2\,m air temperature (Khulna, Chuadanga, Rajshahi; 2024--2025) to capacity, internal resistance, depth of discharge, Arrhenius/Miner-rule ageing, replacement-aware 5- and 10-year discounted cost, and a carbon and lead-flow account. Inputs come from published sources wherever they exist (a 2026 field study of Khulna and Gopalganj easy-bikes, peer-reviewed ageing data, the June 2026 BERC tariff and an official grid emission factor); inputs without a source are run as explicit scenarios. Site-averaged results give a battery life of \lifeLA{} months for lead-acid, \lifeNMC{} for NMC and \lifeLFP{} for LFP, and a 10-year net present cost of BDT \tcoLA{} (lead-acid), \tcoNMC{} (NMC) and \tcoLFP{} (LFP) per bike. LFP cuts grid electricity by \gridratio\% and emits about \saveLFPbike{}~kg\,CO$_2$e less per bike-year than lead-acid ($\approx$\fleetLFPa{}~Mt\,CO$_2$e\,yr$^{-1}$ for 1.4~million bikes), while removing \leadLow--\leadHigh{}~kt of lead per year from the recycling stream. The three sites differ by only \Tspread$^\circ$C in mean battery temperature, so regional temperature does not change the ranking; battery temperature rise, daily distance and energy use matter more, and LFP's cost advantage disappears only when several unfavourable assumptions combine. The study is model-based and has not been validated against operator records.
\end{abstract}
\noindent\textbf{Keywords:} easy-bike; lead-acid; lithium iron phosphate; NASA POWER; battery degradation; total cost of ownership; carbon emissions; lead pollution; Bangladesh
